\documentclass[11pt,letterpaper]{article}
\usepackage[T1]{fontenc}
\usepackage[utf8]{inputenc}
\usepackage{lmodern,microtype}
\usepackage[margin=.92in,headheight=14pt]{geometry}
\usepackage{amsmath,amssymb,amsthm,mathtools}
\usepackage{booktabs,tabularx,array,enumitem}
\usepackage{xcolor,listings,fancyhdr,xurl,needspace}
\usepackage[colorlinks=true,linkcolor=blue!40!black,urlcolor=blue!40!black,citecolor=blue!40!black]{hyperref}
\definecolor{ink}{HTML}{193B4E}
\definecolor{pale}{HTML}{F4F7F8}
\newtheorem{theorem}{Theorem}[section]
\newtheorem{lemma}[theorem]{Lemma}
\newtheorem{proposition}[theorem]{Proposition}
\newtheorem{corollary}[theorem]{Corollary}
\theoremstyle{definition}\newtheorem{definition}[theorem]{Definition}
\theoremstyle{remark}\newtheorem{remark}[theorem]{Remark}
\newcommand{\ipra}{\mathsf{iPRA}}
\newcommand{\Ucal}{\mathcal U^{-}}
\newcommand{\Gcal}{\mathcal G}
\newcommand{\Ncal}{\mathcal N}
\newcommand{\Q}{\mathsf Q}
\newcommand{\D}{\mathsf D}
\newcommand{\FV}{\operatorname{FV}}
\newcommand{\file}[1]{{\small\path{#1}}}
\setlength{\parskip}{.35em}
\setlength{\emergencystretch}{3em}
\setlist{topsep=4pt,itemsep=3pt,leftmargin=1.7em}
\lstset{basicstyle=\ttfamily\footnotesize,columns=fullflexible,
  keepspaces=true,breaklines=true,showstringspaces=false,
  frame=single,rulecolor=\color{ink!20},backgroundcolor=\color{pale},
  aboveskip=9pt,belowskip=9pt}
\pagestyle{fancy}\fancyhf{}
\fancyhead[L]{\small\color{ink}A smaller checked presentation of iPRA}
\fancyhead[R]{\small November 11, 2026}
\fancyfoot[C]{\thepage}
\hypersetup{pdftitle={A Smaller Checked Presentation of iPRA},
 pdfauthor={Research prepared for Vladimir Reshetnikov}}

\begin{document}
\begin{center}
{\LARGE\bfseries\color{ink}A Smaller Checked Presentation of iPRA}\\[7pt]
{\large Immutable proof contexts and a scalar neutral premise}\\[10pt]
Research prepared for Vladimir Reshetnikov\\[3pt]
November 11, 2026
\end{center}

\begin{abstract}
Under ConciseIPRA's complete-proof-checker metric, three standalone Python
modules attain a new coordinate envelope of \textbf{256 significant tokens,
278 AST nodes, and 521 normalized bytes}. The corresponding envelope in
the inspected current primary checker corpus, including all six incoming
archives, is 257 tokens, 279 nodes, and 522 bytes. Different programs attain
the coordinates. The improvement is small and concerns the implementation
of an inherited axiomatization: a global-scope checker is converted to a
loop with immutable contexts, and local-scope checkers use scalar zero as
their neutral true premise. Neither change removes a mathematical axiom.

The implemented public calculus is the existing guarded-source presentation.
It is equiconsistent and bi-interpretable with full first-order iPRA, with
whole numerical domains, actual equality, and identity numerical
composites. We state its entire rule inventory and both interpretations,
prove the preservation properties of the implementation changes, and retain
the inherited mathematical sources. A separate elementary theorem weakens
the scalar kernel's equality rule to replacement of one left-side
occurrence. Its checker-size benefit is not established. No absolute
minimality or new proof-assistant verification is claimed.
\end{abstract}

\section{Result and comparison boundary}
\label{sec:result}

The repository was read at commit
\begin{center}
\texttt{a4e5e96c456303cf231eab2aeeb3291d8889833f}.
\end{center}
The comparison directory is
\file{Logic/PrimitiveRecursiveArithmetic/ConciseIPRA}.
The README describes the native scalar calculus, its Lean equivalence
with iPRA, and twenty-eight filed research packages. The current tree also
contains six incoming ZIP archives. These are part of the comparison,
rather than being treated as future work.\cite{repo,incoming}

The measured improvement is a refinement of
\file{ipra_guarded_source_2026-11-11.zip}, already present in that tree.
The source-dispatch reduction, guarded matrices, and the proof that they
compose are inherited. The scalar-zero neutral-premise convention also
appears in the incoming reversal package; its use here is a composition of
existing techniques.\cite{composition,reversal}

\begin{center}\small
\begin{tabular}{@{}lrrrr@{}}
\toprule
Complete standalone module & Tokens & AST & Bytes & Gzip-9\\
\midrule
Incoming token profile &257&285&568&338\\
New token profile &\textbf{256}&283&594&346\\
Incoming AST profile &261&279&603&345\\
New AST profile &260&\textbf{278}&602&344\\
Incoming byte profile &264&301&522&328\\
New byte profile &263&300&\textbf{521}&327\\
\bottomrule
\end{tabular}
\end{center}

These figures count the \emph{entire} module, including JSON input,
scope handling, substitution, all proof rules, and acceptance output.
The token-oriented change increases its normalized byte length; it is
chosen for the token coordinate. The three coordinate improvements are
approximately $0.389\%$, $0.358\%$, and $0.192\%$, respectively. Gzip is
informational and is not the optimization criterion.

The source audit screened 151 primary Python files under the filed and
incoming \file{code/} directories, retaining 123 standalone checker
candidates. Short generators and support scripts were not counted as
checkers. The detailed rows, exclusions, hashes, and runtime information
are in \file{evidence/primary_checker_audit.json}. This is a documented
comparison corpus, not a lower bound over all possible presentations.

\paragraph{Input contract.}
The compact modules assume that the claimed public theorem is supplied in
the canonical guarded encoding. They may accept additional sound
administrative formulas; they do not validate the entire public grammar.
The separately counted \file{code/checker_strict.py} validates a closed
public goal as well. It is a \emph{public-goal} validator, not a claim that
every object in the certificate has undergone a separate exact-type check.
Its complete size is 1,201 normalized bytes (489 tokens, 560 AST nodes,
571 gzip-9 bytes), versus 1,202 bytes for the identical inherited validator
with its old neutral premise. Its metrics and tests are recorded alongside
the compact profiles. Comparisons to validators for other public grammars
would be comparisons under different contracts.

\section{The complete implemented axiomatization}
\label{sec:calculus}

This section fixes the mathematical presentation implemented by the new
modules. It is the inherited guarded-source calculus $\Ucal$ with its
specified public language. All primitive axioms and rules are given here.

\subsection{Terms and formulas}

There is one numerical sort. Terms are finite expressions
\[
 t ::= x\mid 0\mid S(t)\mid I(n,a;u.h).
\]
Only $u$ in $h$ is bound. Other variables in $h$ are fixed numerical
parameters. Equivalently, each fixed finite earlier-stage body $h$ gives a
first-order function symbol $I_h$. No numerical variable ranges over
function descriptions, and there is no uniformly total universal evaluator.
Write $1=S0$ and
\[
 C(v,a,b):=I(v,a;u.b).
\]

Public formulas are equations and paired binders $\D_{x,y}B$.
An equation immediately under $\D$ must have the exact shape
\begin{equation}
 E(x,y)\equiv I(y,0;u.h)=I(y,0;u.k).
 \tag{Guard}\label{eq:guard}
\end{equation}
The update bodies may depend on $x,y,u$ and passive parameters. If the
matrix is compound, its root is another $\D$, and it must itself be public.
Every guarded equation has a provable zero instance. The designated
contradiction is the public equation $0=1$.

Proofs may additionally use the auxiliary formulas
\[
 A ::= (s=t)\mid \Q_y A\mid \D_{x,y}A,
\]
including unguarded atomic matrices. A public sequent has public assumptions
and conclusion; its internal proof formulas need not be public. This is an
explicit feature of the presentation, and the equivalence theorem below
does not assert both round trips for every auxiliary formula.

\subsection{All primitive axioms and rules}

The only numerical axiom schemas, available in every retained context, are
\begin{align}
 I(0,a;u.h)&=a,\tag{I0}\label{eq:i0}\\
 I(Sn,a;u.h)&=h[I(n,a;u.h)/u].\tag{IS}\label{eq:is}
\end{align}
Substitution is capture-avoiding throughout. There are four additional
inference patterns:

\paragraph{Hypothesis.}
\[
 \frac{}{\Gamma\vdash A}\quad(A\in\Gamma).
 \tag{Hyp}
\]

\paragraph{Source-selected elimination.}
For an auxiliary formula template $F(v)$, define
\[
 M(s,t;F)=
 \begin{cases}
  \Q_vF(v),&s\text{ is literally }0,\\
  s=t,&\text{otherwise}.
 \end{cases}
\]
Then use
\[
 \frac{\Gamma\vdash M(s,t;F)\qquad\Gamma\vdash F(s)}
      {\Gamma\vdash F(t)}.
 \tag{SE}\label{rule:se}
\]
The restriction is syntactic: a nonzero-looking iterator term that denotes
zero is an allowed equality source. At literal zero, this rule is anchored
closure elimination.

\paragraph{Paired elimination.}
\[
 \frac{\Gamma\vdash\D_{x,y}B(x,y)\qquad
       \Gamma,\Q_yB(b,y)\vdash A}
      {\Gamma\vdash A}.
 \tag{DE}\label{rule:de}
\]
Here $b$ is fresh for $\Gamma$, $A$, and the passive parameters of the
matrix. The displayed local assumption is discharged.

\paragraph{Deterministic paired introduction.}
Put $\eta_B(z)=z$ when $B$ is a literal equation, and $\eta_B(z)=0$
otherwise. Use
\[
 \frac{\Gamma,B(t,\eta_B(z))\vdash B(t,Sz)}
      {\Gamma\vdash\D_{x,y}B(x,y)}.
 \tag{DI}\label{rule:di}
\]
The eigenparameter $z$ is fresh for the retained context, the witness $t$,
and passive matrix parameters. Again, the local assumption is discharged.

This inventory is complete. Reflexivity, symmetry, numerical coverage,
successor injectivity and separation, equation induction, ordinary
connectives, and ordinary quantifiers are derived. In particular, no rule
of classical excluded middle or induction on arbitrary formulas is added.

Weakening, assumption cut, and fresh renaming are admissible. Numerical
substitution into a restricted-source proof is not used as a primitive
metarule: replacing a variable source by literal zero can invalidate a
primitive SE step. Later equivalence theorems may establish substitution
as a derived operation on the public consequence relation.

\section{Why the presentation is exactly iPRA}
\label{sec:equivalence}

We give the interpretations and the dependency order, rather than relying
only on the assertion that the standard numerical functions are primitive
recursive. The detailed inherited proofs are compiled in
\file{foundations.pdf}, with their source modules retained unchanged.

\subsection{Auxiliary soundness and public interpretation}

Interpret every finite iterator description by its corresponding primitive
recursive definition. Preserve numerical equality. Auxiliary formulas have
the syntax-sensitive reading
\begin{align*}
 (s=t)^*&=(s=t),\\
 (\Q_yB)^*&=B^*(0)\to\forall y\,B^*(y),\\
 (\D_{x,y}B)^*&=
 \begin{cases}
 \exists x\,\forall y\,(B(x,y)\to B(x,Sy)),&B\text{ an equation},\\
 \exists x\,(B^*(x,0)\to\forall y\,B^*(x,y)),&B\text{ compound}.
 \end{cases}
\end{align*}
Every inference is sound in iPRA. Atomic DI introduces a progressive
predicate. Atomic DE uses equation induction to obtain its anchored
closure. Compound DI uses the zero-or-successor coverage theorem, and
compound DE is existential elimination. SE is equality transport or anchored
closure elimination. Formula roots are preserved by numerical substitution,
so this translation commutes with that substitution.

The public reading instead interprets every $\D$ uniformly by
\begin{equation}
 (\D_{x,y}B)^\circ
   =\exists x\bigl(B^\circ(x,0)\to\forall y\,B^\circ(x,y)\bigr).
 \label{eq:ordinary}
\end{equation}
At a guarded equation, $E(x,0)$ is a theorem. Hence iPRA proves
\[
 \forall y(E(x,y)\to E(x,Sy))
 \ \longleftrightarrow\
 E(x,0)\to\forall yE(x,y).
\]
The forward direction uses only equation induction; the reverse uses the
true base instance. Structural induction therefore proves
$A^*\leftrightarrow A^\circ$ for every public $A$.

\subsection{Completeness without a circular substitution argument}

The inherited completeness proof proceeds in the following order.

\begin{enumerate}
\item Derive contextual equation induction in the auxiliary calculus.
From $E(z)\vdash E(Sz)$, atomic DI with an unused witness gives
$\D_{x,y}E(y)$. DE exports $\Q_yE(y)$, and SE applies it to the
base $E(0)$ at any desired numerical target. The equation may be unguarded.

\item Work first in the purely equational fragment. Computation and
transport with sources other than literal zero, together with the derived
equation induction, suffice to reverse $0=t$ into $t=0$. The supplied proof
uses a represented zero, a singleton test, and an unconditional zero filter.
Equation cases come directly from equation induction. The proof keeps
passive parameters fixed and does not substitute zero into a derivation
whose transport source was a variable.

\item Recover unrestricted numerical equality reasoning and then
transport through every public formula template. At a guarded atomic
matrix its zero instance is true, so the opened closure supplies the
successor instance needed for DI. At a compound matrix, use induction on
the template to transport the local base backwards, apply the opened
closure, and transport the successor instance forwards.

\item Obtain exactly the consequences of the inherited guarded calculus
$\Gcal$ on public contexts and conclusions. In the reverse direction,
translate a proof into iPRA using the mixed reading, replace it by the
equivalent public reading, and apply the faithful guarded interpretation
back. No identity claim for an arbitrary auxiliary formula is used.
\end{enumerate}

The distinction in the last step matters. The administrative syntax serves
as a sound proof language; the public syntax carries the bi-interpretation.
The supplied mathematical review checks this dependency order explicitly.

\subsection{Both translations and their composites}

The ordinary scalar theory $T_{\rm eq}$ uses the iterator equations,
intuitionistic first-order equality logic, numerical coverage,
$0=1\to\bot$, and contextual induction for literal equations. Its inherited
arithmetic construction derives predecessor, equality reflection into a
Boolean term, quantifier-free induction, pairing and projections, and then
a compiler for every fixed primitive-recursive description.\cite{foundations,lean}

For example, after pairing is available, a primitive recursion
\[
 f(0,\vec p)=g(\vec p),\qquad
 f(Sn,\vec p)=h(n,f(n,\vec p),\vec p)
\]
is compiled by iterating the state update
\[
 r\longmapsto
 \langle S\pi_0r,\ h(\pi_0r,\pi_1r,\vec p)\rangle
\]
from $\langle0,g(\vec p)\rangle$, and projecting the second coordinate.
The counter and value use one scalar state. The proof of the defining
equations and of each round trip is internal equation induction, rather
than an external observation about standard numerals.

For the logical translation, let $\operatorname{sel}(v,A,B)$ select $A$
at zero and $B$ at successors, and let $\operatorname{gate}(v,A)$ be true
at zero and equivalent to $A$ at successors. Their construction uses
the numerical selector $C$ at equational leaves and finite syntactic
padding of compound formulas. In the native Eq/$\Q$/$\D$ calculus,
\[
\begin{aligned}
 \bot&:=0=1,\\
 A\to B&:=\Q_y\operatorname{sel}(y,A,B),\\
 A\wedge B&:=\Q_y\operatorname{gate}
      (y,\operatorname{sel}(\operatorname{par}(y),A,B)),\\
 \exists x A(x)&:=\D_{x,y}\operatorname{gate}(y,A(x)),\\
 A\vee B&:=\exists v\,\operatorname{sel}(v,A,B),\\
 \forall x A(x)&:=A(0)\wedge\Q_y A(y).
\end{aligned}
\]
Here $\operatorname{par}$ is a fixed parity iterator. The inherited proofs
derive all natural-deduction rules for these macros.

Replace each $\Q_yB$ by $\D_{u,y}B$ with an unused fresh witness $u$.
To make matrices public, for an equation $E\equiv s=t$ put
\[
 L(E):=\D_{u,v}\bigl(C(v,0,s)=C(v,0,t)\bigr),
 \qquad
 P(A):=\begin{cases}L(A),&A\text{ an equation},\\A,&A\text{ compound}.
 \end{cases}
\]
The padding translation is
\[
 T(s=t)=s=t,\qquad T(\D_{x,y}B)=\D_{x,y}P(TB).
\]
Both padding composites are internally interderivable with the original
formula; the guarded atomic case uses its true zero instance.

Let $\Phi$ be the iPRA-to-public translation obtained by composing the
scalar compiler, logical macros, dummy-witness expansion, and padding.
Let $\Psi$ be the public ordinary reading followed by expansion of each
scalar iterator into its primitive-recursive definition. Then
\begin{align*}
 \ipra&\vdash\varphi\leftrightarrow\Psi\Phi(\varphi),\\
 A&\vdash_{\Ucal}\Phi\Psi(A)
       \quad\text{and}\quad\Phi\Psi(A)\vdash_{\Ucal}A.
\end{align*}
The second line holds for every public $A$. The interpretations
use the whole numerical domains and actual equality, and the graph of each
numerical composite is $x=y$. Consequently
\[
 \Ucal\vdash 0=1\quad\Longleftrightarrow\quad\ipra\vdash\bot.
\]
These are the same interpretations for every new checker: the
implementation changes do not change the public consequence relation.

\section{The implementation reductions and their proofs}
\label{sec:implementation}

\subsection{Certificate conventions}

Each module reads a JSON array from standard input. Its last item is the
goal; preceding instruction packets are consumed from right to left. A
packet is $[m,s,u,t]$. The checker either closes a retained hypothesis,
uses a selected elimination premise, supplies a computation equation, or
introduces a paired formula. An unused prefix is inert. Acceptance means
normal process exit with exactly \texttt{ACCEPT} and a newline. Exceptions,
nonzero exit, and resource failure are not acceptance. Assertions are not
logical gates.

The inherited arity codec uses unary arrays for $\D$ formulas, binary
arrays for equations, and ternary arrays for administrative $\Q$ closures.
Numerical zero is scalar $0$, successor is unary, and the iterator is
ternary. The opening traversal increments its depth by
$5-\operatorname{arity}$; numerical roots use depth $3$. The encoding's
protected and dummy coordinate conventions are specified in the retained
arity source and certificate documentation. Extra raw shapes have an
administrative interpretation; they are not silently new public formulas.

There are two protocols. Local profiles carry the fresh-variable floor
as an argument of each call. The token profile uses a shared global floor.
They do not in general accept the same serialized certificate. The archive
labels and tests the two sets separately.

\subsection{Tail-call elimination with immutable contexts}

The incoming token checker ends a non-hypothesis step by recursively checking
the continuation with an extended context. Replace that tail call by a loop
and use a tuple for the retained context:
\begin{lstlisting}[language=Python]
while i(a) not in g:
    # Check the instruction and its major premise.
    ...
    k -= 1
    g += (i(h, (m or len(a) == 2) and k),)
    a = i(a, [k])
\end{lstlisting}
The complete measured source is \file{code/checker_tokens.py}; this
illustration is not the counted program.

\begin{theorem}[Continuation preservation]
Under ideal unbounded finite-certificate execution, replacing the incoming
token checker's final recursive call by this loop and replacing contexts by
tuples preserves its successful proof executions, packet consumption, and
global fresh-name state.
\end{theorem}
\begin{proof}
Relate an old list context and a new tuple context when they have the same
entries in the same order. Membership tests agree. A major-premise call
receives related contexts and the same packet stream and global floor;
by induction on its finite successful call tree it returns after consuming
the same packets with the same floor. The next decrement is therefore the
same in both executions.

The old final call receives the opened goal and the old context concatenated
with one opened hypothesis. The loop next checks precisely those objects
with the same global floor. Tuple concatenation creates a new context value
and rebinds only the local name. It does not change the context retained by
the caller of a major-premise proof. The pure opening operations are
evaluated in the opposite order in the two continuation spellings, but
neither changes the packet stream, context, or global floor. On successful
finite executions both finish and give identical objects.

When a hypothesis closes, both versions return \texttt{None}; the old tail
frames then only return \texttt{None} as well. This gives both simulation
directions on finite successful executions.
\end{proof}

Fixed CPython stack limits need not give identical resource failures: the
loop can check some longer valid continuations. The theorem compares proof
recognition under the same ideal execution convention as the baseline.

\paragraph{Why a tuple is necessary.}
Using a list with \texttt{g += [hypothesis]} would mutate a context shared
with the caller and permit a local assumption to escape. The included
regression certificate
\begin{lstlisting}[language={}]
[[[],0,[],[0]],[[[0,[0]]],0,0,0],[0,[0]]]
\end{lstlisting}
claims the false public equation $0=1$. The correct tuple checker rejects
it. The deliberately constructed mutable-list mutant accepts it. This is a
test of a logical scope property, not merely a comparison of syntax counts.

\subsection{Scalar neutral premise}

The local profiles initially retain a true padding formula. Replace its
spelling \texttt{[]} by scalar \texttt{0}, and use \texttt{0} in the neutral
major field of local certificates:
\begin{lstlisting}[language=Python]
def c(a, g=[0], k=0):
    ...
\end{lstlisting}
The tuple-context AST profile analogously uses \texttt{g=(0,)}.

\begin{proposition}[Neutral-premise preservation]
The scalar-neutral local profiles remain sound for their inherited raw
administrative reading and complete for the same public guarded
consequences.
\end{proposition}
\begin{proof}
Both neutral spellings mean truth at formula sort. For a neutral major,
the elimination match fails in either spelling, the recursive premise check
closes immediately, its tripled export is false as a Python condition, and
the fresh-assumption selector takes the same branch. Therefore a canonical
proof certificate is translated by replacing neutral-major fields only;
other fields and the public goal remain unchanged.

In canonical derivation certificates every meaningful formula goal has
nonempty root arity, which opening preserves. It cannot be confused with
the new scalar neutral assumption. This proves completeness by translating
the existing certificate construction. For arbitrary raw certificates,
the only added initial resource is still a true formula, and all actual
rule checks are unchanged. The inherited raw soundness argument therefore
applies with the new true resource.
\end{proof}

This proposition is not a claim that the two modules accept exactly the
same raw JSON language. Their neutral premise spelling differs. The
translation acts on packet fields, not on every empty array occurring
inside a formula or term.

\paragraph{A restricted-substitution detail.}
The arity checker's computation-supply proof admits fresh numerical
parameters and subsequently specializes them to zero. In this restricted
calculus, make that specialization \emph{after} interpreting the continuation
in ordinary iPRA, where numerical substitution is valid. Then cut the known
computation instance. Do not apply an unrestricted substitution theorem
directly to a restricted-source derivation. The inherited guarded-source
proof uses exactly this order.

\section{An auxiliary mathematical weakening of equality}
\label{sec:lor}

The following result concerns the unrestricted-source native scalar kernel
$\Ncal$ formalized in the repository. It is independent of the measured
implementation improvement and does not replace the implemented SE rule.

\begin{definition}[Left-occurrence replacement]
Retain all rules of $\Ncal$ except equational transport, and replace that
rule by
\[
 \frac{\Gamma\vdash s=t\qquad\Gamma\vdash C[s]=r}
      {\Gamma\vdash C[t]=r},
 \tag{LOR}
\]
where $C[-]$ is a capture-avoiding numerical term context with exactly one
hole; the identity context is allowed. The right endpoint $r$ is fixed.
\end{definition}

\begin{theorem}[One occurrence suffices]
For every native context and formula, the LOR calculus and $\Ncal$ derive
exactly the same sequents. Their term and formula translations are identities.
\end{theorem}
\begin{proof}
Write $j_a=I(0,a;u.0)$. Computation gives $j_a=a$. Use this equation as
both premises of LOR with the identity context to obtain $a=a$.
From $s=t$ and $s=s$, LOR gives $t=s$.

Given an old transport step with template $A(v)=B(v)$ and premises
$s=t$, $A(s)=B(s)$, enumerate the distinguished free occurrences of $v$
in $A$, and replace them one at a time by LOR. This gives $A(t)=B(s)$.
Apply the just-derived symmetry, replace the distinguished occurrences in
$B$ one at a time, and symmetrize again. The result is $A(t)=B(t)$.
Before replacement, rename bound iterator variables fresh from $s,t$ and
the passive parameters. Occurrences are tracked in the template, not found
by replacing every subtree that happens to equal $s$.

Conversely, LOR is an instance of the old rule with equation template
$C[v]=r$. Replacing each old transport node by the finite construction
above translates derivations in either direction, in every retained context.
\end{proof}

The repository's native/iPRA bi-interpretation transfers immediately to
this alternative kernel. The primitive equality interface is more
restricted, but it has not been shown to yield a smaller complete checker:
the old checker shares formula opening between transport and closure
elimination. A new occurrence-position traversal has a real cost. This is
elementary equality reasoning, and the earlier counter-congruence package
already studies related restrictions in another numerical language.

\section{A tempting reduction introduces full induction}
\label{sec:barrier}

The root test distinguishing equation progress from compound cases is a
logical boundary. An attempted implementation reduction uses the generated
hypothesis's arity instead of both the major mode and the matrix root.
Administrative closures have arity three, so that reduction also permits
progress at $\Q$-rooted matrices.

\begin{proposition}[Closure progress yields unrestricted induction]
With the ordinary native reading of $\Q$, permitting progress introduction
at arbitrary $\Q$-rooted matrices yields induction for every native formula.
\end{proposition}
\begin{proof}
For an arbitrary $P(n)$ define
\[
 G(n):=\Q_z\operatorname{gate}(z,P(n)).
\]
The gate is true at zero and equivalent to $P(n)$ at each successor, so
$G(n)\dashv\vdash P(n)$. A proof of $P(n)\vdash P(Sn)$ therefore gives
$G(n)\vdash G(Sn)$. The proposed progress introduction, with a dummy
witness, followed by DE and QE supplies induction for $G$. Its base and
conclusion translate respectively to $P(0)$ and $P(t)$. Since $P$ was
arbitrary, the result is full formula induction.
\end{proof}

The strengthening reaches public Eq/$\D$ formulas through the existing
translation. Restricting only public goal syntax therefore does not repair
it. One can instead distinguish object formulas and closure resources by a
genuine type discipline, or prove a normalization theorem allowing only
hypothesis checks at closure goals. Those costs must be included in the
checker; merely giving the same raw object different meanings in different
uses fails to justify closure elimination.

\section{Verification, reproduction, and further questions}
\label{sec:verification}

The archive contains the three compact modules, an expanded standalone
checker, a public-goal validator, frozen comparison modules, the original
global and local certificates and their translated forms, and a standard-
library reproduction script. The script measures complete modules and runs
positive and negative controls normally and with Python optimization.
It also executes the context-leak regression against the correct checker
and a temporary mutant.

The final reproduction run completed \textbf{210 candidate/control comparisons} (420 Python processes) and four additional scope-control executions: \textbf{424 processes in total, with zero discrepancies from their prescribed outcomes}. The correct checker rejects the scope attack in both modes; the deliberate mutable-list mutant accepts it in both.


The checks include computations, reflexivity, a used equation-induction
hypothesis, compound case introduction, zero-source transport, malformed
packets, false equations, scope leakage, and future-witness capture. The
receipts list every tested process and verdict. These finite executions
support the implementation changes; they do not replace the universal
proof arguments or constitute a Lean/Rocq verification.

The existing Lean native/iPRA bridge was read at the pinned commit. It was
not rebuilt here. The guarded-source bridge is inherited written
mathematics, independently reviewed in this run. The review notes state
their scope explicitly; they are not external human peer review.

\paragraph{Reproduce.}
From the extracted archive root, run
\begin{lstlisting}[language=bash]
python3 reproduce.py
pdflatex -interaction=nonstopmode -halt-on-error report.tex
pdflatex -interaction=nonstopmode -halt-on-error report.tex
pdflatex -interaction=nonstopmode -halt-on-error foundations.tex
pdflatex -interaction=nonstopmode -halt-on-error foundations.tex
pdflatex -interaction=nonstopmode -halt-on-error foundations.tex
\end{lstlisting}
The checker and reproduction commands require Python's standard library.
The PDF commands require the ordinary LaTeX packages named in the sources.
The archive records CPython and zlib versions because normalized source
and compression measurements can depend on them.

\subsection*{Further research questions}

\begin{enumerate}
\item Can one complete module beat all three current coordinate minima
simultaneously, without changing the canonical-goal contract? The present
result uses separate programs.
\item Can public syntax validation share the kernel's opening traversal
while still checking its binding conventions completely? This is a more
substantial engineering target than another one-byte permissive reduction.
\item Can the guarded-source dependency theorem and both public composites
be added to the existing Lean formalization, with the restricted
substitution boundary explicit?
\item Can closure resources be given an intrinsic representation that
permits a uniform continuation rule at lower total cost? A proof must
include all eliminations, not only the introduction rule.
\item Does left-occurrence transport help a readable or formally verified
checker even when it does not help the compact arity implementation?
\item Is there a different arithmetic foundation that removes an operation
and also reduces the complete checker? This investigation has not
established such a mathematical improvement.
\end{enumerate}

\begin{thebibliography}{9}
\bibitem{repo} Vladimir Reshetnikov and ProveIt contributors,
\emph{Concise presentations of intuitionistic primitive recursive arithmetic},
pinned README, commit \texttt{a4e5e96c456303cf231eab2aeeb3291d8889833f}.
\url{https://github.com/VladimirReshetnikov/ProveIt/blob/a4e5e96c456303cf231eab2aeeb3291d8889833f/Logic/PrimitiveRecursiveArithmetic/ConciseIPRA/README.md}.
\bibitem{incoming} ProveIt, \emph{ConciseIPRA incoming proposals}, same commit.
All six ZIP archives were inspected; exact member provenance and primary
checker measurements are supplied with this report.
\url{https://github.com/VladimirReshetnikov/ProveIt/tree/a4e5e96c456303cf231eab2aeeb3291d8889833f/Logic/PrimitiveRecursiveArithmetic/ConciseIPRA/docs/incoming}.
\bibitem{composition} ProveIt research package,
\emph{A smaller guarded presentation of iPRA}, member \file{report.tex}
of \file{ipra_guarded_source_2026-11-11.zip} in the preceding directory.
Source of the unchanged implemented calculus, mixed auxiliary reading,
equational dependency theorem, public consequence theorem, and three
comparison modules.
\bibitem{reversal} ProveIt research package,
\emph{A smaller checked presentation of iPRA},
\file{ipra_reversal_core_2026-11-11.zip}, same incoming directory.
Source of the previously used scalar-zero neutral-premise convention.
\bibitem{foundations} ProveIt,
\emph{certified\_zero/reference}, \emph{guarded\_core}, and
\emph{source\_dispatch}, same pinned commit. Attributed frozen mathematical
sources are in \file{reference/foundations/}; the supplementary
\file{foundations.pdf} compiles the required mathematical chain.
\url{https://github.com/VladimirReshetnikov/ProveIt/tree/a4e5e96c456303cf231eab2aeeb3291d8889833f/Logic/PrimitiveRecursiveArithmetic/ConciseIPRA/Research}.
\bibitem{lean} ProveIt contributors, \emph{ConciseIPRA Lean development},
especially \file{NativeBiInterpretation.lean}, \file{BiInterpretation.lean},
and \file{NativeZeroSource.lean}, same commit.
\url{https://github.com/VladimirReshetnikov/ProveIt/tree/a4e5e96c456303cf231eab2aeeb3291d8889833f/Logic/PrimitiveRecursiveArithmetic/ConciseIPRA/Lean/ConciseIPRA}.
\end{thebibliography}
\end{document}
